% Section 1.5, from lectures/L01/README.md: what you should be able to do, questions to test
% yourself, the next lecture, and the references.

\section{Review}
\label{c1:sec:review}

\needspace{6\baselineskip}
\subsection*{What you should be able to do now}
\begin{itemize}
  \item Say what the 32 registers are, which sixteen an immediate instruction can name, and why
    the other sixteen cannot be named that way.
  \item Read a line of \code{SREG}'s contents and say which branch would be taken next.
  \item Given an address, say which of the three address spaces it belongs to and which
    instruction reaches it; and say why \code{in}, \code{lds} and \code{lpm} are three different
    instructions rather than one.
  \item Encode an \code{ldi} by hand and confirm it against a disassembly.
  \item Assemble a \code{.asm} file, disassemble the result, and run it in the simulator.
  \item Count the cycles of a straight-line routine and of a loop, from the instruction set
    summary, and say why a measured figure at a call site comes out higher than your count of the
    body.
  \item Write a subroutine that takes an argument in \code{r24}, returns a value in \code{r24},
    and returns.
  \item Write a program that runs on its own: a reset vector, a stack pointer, an LED lit through
    \code{DDRB} and \code{PORTB}, and a loop to stop in.
\end{itemize}

\needspace{6\baselineskip}
\subsection*{Questions to test yourself}
\begin{enumerate}
  \item Why can \code{ldi} name \code{r16} but not \code{r15}, and what does that cost you when you
    run out of registers in the upper half?
  \item Data space address \code{0x0025} and I/O address \code{0x05} are the same register. Which
    instructions take which of the two numbers, and what happens if you give one the other's?
  \item The register file is addressable as data at \code{0x0000} to \code{0x001F}. What could you
    do with that, and why does essentially no program do it?
  \item A conditional branch costs one cycle when it falls through and two when it branches. A
    loop running $n$ times executes its exit branch $n + 1$ times. How many of those are taken?
  \item Your hand count of a subroutine says 40 cycles and the simulator agrees, yet the
    \emph{call} costs 44 and the loop around it costs 46 per iteration. Where are the other four,
    and the two after that, without looking anything up?
  \item Why does \code{shift_bits(7)} cost more than five times what \code{shift_bits(0)} costs,
    and what does that imply about a driver that calls it with a pin number?
  \item Your program writes \code{DDRB} at I/O address \code{0x04}, and the test that checks it
    reads data space address \code{0x24}. Both are right. Which instructions take which of the two
    numbers?
  \item You light an LED with \code{sbi PORTB, 5} and it works. Why can a driver that has to light
    \emph{any} LED not be written that way at all?
\end{enumerate}

\needspace{6\baselineskip}
\subsection*{Looking ahead}
Chapter~2 takes the same LED and lights it with something that can be told which LED to light:
\begin{itemize}
  \item \code{DDRx}, \code{PORTx} and \code{PINx}: three registers per port, and what each of the
    three actually does, which is more than the two \code{sbi} instructions at the end of this
    chapter needed to know.
  \item Why writing a one to \code{PINx} toggles an output, which is not what the name suggests at
    all.
  \item The calling contract: which registers a subroutine may destroy, and which it must give
    back.
  \item The LED driver: the same LED, lit by something that can be told which LED to light.
\end{itemize}

\needspace{6\baselineskip}
\subsection*{Sources}
Every number this chapter quotes comes from one of two primary documents, and the text says
which:
\begin{itemize}
  \item \emph{ATmega328P Datasheet}, Microchip Technology.\\
    \url{https://ww1.microchip.com/downloads/en/DeviceDoc/Atmel-7810-Automotive-Microcontrollers-ATmega328P_Datasheet.pdf}
  \item \emph{AVR Instruction Set Manual}, Microchip Technology.\\
    \url{https://ww1.microchip.com/downloads/en/devicedoc/atmel-0856-avr-instruction-set-manual.pdf}
\end{itemize}
